\documentclass[10pt,a4paper]{amsart}
\usepackage[margin=19mm]{geometry}
\usepackage{amsmath,amssymb,mathtools}
\usepackage[scr=rsfs,cal=euler]{mathalfa}
\usepackage{xfrac}
\usepackage[unicode,hidelinks,bookmarks=false]{hyperref}
\newcommand{\ie}{i.e.\ }
\newcommand{\cf}{cf.\ }
\newcommand{\resp}{respectively\ }
\newcommand{\Subsection}{\subsection}
\providecommand{\nobreakdash}{\nobreak\hbox{-}}
\setlength{\emergencystretch}{2em}
\multlinegap0pt
\usepackage{siunitx}
\usepackage{siunitx}
\usepackage{siunitx}
\usepackage{siunitx}
\usepackage{amssymb}
\usepackage{enumerate}
\usepackage{xcolor}
\usepackage{dsfont}
\usepackage{tikz}
\usepackage{siunitx}
\newtheorem{defn}{Definition}
\newtheorem{lem}{Lemma}
\newtheorem{thm}{Theorem}
\usepackage{cleveref}
\crefname{defn}{Definition}{Definitions}
\crefname{lem}{Lemma}{Lemmas}
\crefname{thm}{Theorem}{Theorems}
\def\R{{\Bbb R}}
\newcommand{\di}[1]{\,\mathrm{d}#1}
\title{A mathematical model for colloids deposition in porous media combined with a moving boundary at the microscale: Solvability and numerical simulation}
\author{Christos Nikolopoulos, Michael Eden,, Adrian Muntean}
\date{Source: arXiv:2604.09637v1 (CC BY 4.0)}
\begin{document}
\maketitle

\noindent\emph{Source and licence.} A mathematical model for colloids deposition in porous media combined with a moving boundary at the microscale: Solvability and numerical simulation. Version arXiv:2604.09637v1 (CC BY 4.0), distributed by arXiv under the Creative Commons Attribution 4.0 International licence, http://creativecommons.org/licenses/by/4.0/, as recorded in the arXiv licence metadata for this exact version and shown on the abstract page https://arxiv.org/abs/2604.09637v1. Full licence notice: \texttt{licenses/arxiv-2604.09637-CC-BY-4.0.txt}.\par\smallskip

\noindent\textit{Editorial scope.} Counted unit: the article's thm thm.uniqueness (source0.tex lines 1010--1013). The article's complete original proof of it is delivered with it (source0.tex lines 1014--1148, one \texttt{proof} environment of 135 lines). No mathematics is written, repaired or re-derived: the statement and the proof are the article's own lines, in the article's own notation, and no retained line is substituted or re-flowed.  Retained uncounted as the source-local context the proof needs: lines 806--840; lines 886--892; lines 917--925; lines 948--954.  Nothing was omitted from any retained span, so the package carries no omission record.  No retained line contains a citation, so the wrapper carries no bibliography and no bibliography entry.  No retained line contains a figure environment, an image-inclusion command or drawing code, so no image is delivered and the package declares no asset dependency.  Editorial formatting only, in addition: an AMS \texttt{amsart} preamble that restates the article's own declarations and macros used by the retained lines, namely the environments \texttt{defn}, \texttt{lem}, \texttt{thm} on separate counters, together with the standard packages those lines need.  The retained environments print in this document as Definition 1, Lemma 1, Theorem 1 in order of appearance, whereas the article prints the same items with the numbering of its own full document.  The complete native source of the article as distributed in the arXiv e-print for this exact version is bundled unchanged as \texttt{sources/198220/source0.tex}.
\begin{defn}[Weak solutions]\label{def:weak_solution}
A weak solution is a set of functions $(u,v,\sigma)$ with the regularity
%
\begin{align*}
u\in \mathcal{W}(S_T)\cap L^\infty(\Omega\times S_N))^N,\quad \sigma\in L^\infty(\Omega;W^{1,\infty}(S_T)),\quad
v\in L^\infty(\Omega;W^{1,\infty}(S_T))
\end{align*}
that satisfies the following conditions:
\begin{itemize}
    \item[$(i)$] It holds $-\sigma^*\le\sigma(x,t)\le\sigma^*$.
    \item[$(ii)$] The diffusivity of the $i$-nomeres $D_i\in L^\infty(\Omega\times S_T)^{d\times d}$ is given by ($j,k=1,...,d$)
    \[
    (D_i(x,\sigma(t,x))_{jk}=d_i\phi(x,\sigma(t,x))\int_{Y^f(x,\sigma(t,x))}(\nabla w_{k}(x,\sigma(x,t),y)+e_{k})\cdot e_j\di{z}.
    \]
    Here, for any $\sigma\in[-\sigma^*,\sigma^*]$ and $k=1,...,d$ the cell functions $w_{k}(x,\sigma,\cdot)\in H^1_\#(Y(x,\sigma))$ are the unique zero-average solutions to the weak form
    %
    \[
    \int_{Y^f(x,\sigma)}\nabla w_k(x,\sigma,y)\cdot\nabla\xi(y)\di{y}=\int_{\Gamma(x,\sigma)}\xi(y)e_k\cdot n\di{\gamma}
    \]
    for all $\xi\in H^1_\#(Y(x,\sigma))$.
    \item[$(iii)$] The weak forms
    \begin{align*}
    \langle\partial_t u_i,\phi\rangle_{H^1(\Omega)^*}-(D_i(\sigma)\nabla u_i,\nabla\phi)_{L^2(\Omega)} &=(R_i(u),\phi)_{L^2(\Omega)}-(A(\sigma)(a_iu_i- b_iv),\phi)_{L^2(\Omega)}
    \end{align*}
    %
    hold true for all test functions $\phi\in H^1(\Omega)$ and all $1\leq i\leq N$.
    \item[$(iv)$] The equalities
    \begin{align*}
     \partial_tv=\sum_{i=1}^N(a_iu_i- b_iv),\qquad
     \partial_t\sigma=\alpha_v\sum_{i=1}^N(a_iu_i- b_iv).
    \end{align*}
    hold almost everywhere in $\Omega\times S_T$.
\end{itemize}
%
\end{defn}

\begin{lem}[Lipschitz estimates for geometric coefficients]\label{lemma:lipschitz_coefficient}
    There is a Lipschitz constant $C>0$ such that
    \[
    |A(x,\sigma_1)-A(x,\sigma_2)|\leq C|\sigma_1-\sigma_2|
    \]
    for all $x\in\Omega$ and $\sigma_1,\sigma_2\in[-\sigma^*,\sigma^*]$.
\end{lem}

\begin{lem}\label{lemma:lipschitz.diffusivity}
There is a constant $c_D>0$ such that
%
$
c_D|\xi|^2\leq D_i(x,\sigma)\xi\cdot\xi
$
for all $\xi\in\R^d$, for all $x\in\R^n$, and for all $\sigma\in[-\sigma^*,\sigma^*]$.
Moreover, it holds $D_i\in L^\infty(\Omega;C^{1,1}([-\sigma^*,\sigma^*];\R^{d\times d}))$.
\end{lem}

\begin{thm}[Existence]\label{thm.existence}
There is a time interval $S_T=(0,T)$ and a local-in-time weak solution 
\[
(u,v,\sigma)\in \mathcal{W}(S_T)\cap L^\infty(\Omega\times S_T)\times L^\infty(\Omega;W^{1,\infty}(S_T))\times L^\infty(\Omega;W^{1,\infty}(S_T))
\]
solving the Problem in the sense of \cref{def:weak_solution}.
\end{thm}

\begin{thm}[Weak-strong uniqueness]\label{thm.uniqueness}
    If there is a weak solution $(u,v,w,\sigma)$ with the additional regularity $u\in L^2(S_T;W^{1,p}(\Omega))^N$ for some $p>d$.
    Then, this is the unique weak solution.
\end{thm}

\begin{proof}
We assume that we have two sets of solutions $(u^{(j)},v^{(j)},w^{(j)},\sigma^{(j)})$ ($j=1,2$) and denote their difference by $(\bar{u},\bar{v},\bar{w}_k,\bar{\sigma})$.
The surface concentrations are given as
%
\[
v^j(t,x)=e^{-bt}\left(v_0(x)+\sum_{i=1}^Na_i\int_0^te^{b\tau}u_i^j(\tau,x)\di{\tau}\right)
\]
and they are related to the offset parameter $\sigma^{(j)}$ via $\sigma^{(j)}=\alpha_v (v^j-v_0)$.
Therefore, 
\[
\bar{\sigma}(t,x):=\sigma^{(1)}(t,x)-\sigma^{(2)}(t,x)=\alpha_v\sum_{i=1}^N a_i\int_0^te^{b(\tau-t)}\bar{u}_i(\tau,x)\di{\tau}
\]
and
%
\begin{equation}\label{eq:bar_sigma_estimate}
|\bar{\sigma}(t,x)|\leq\alpha_v\sum_{i=1}^N a_i\int_0^t|\bar{u}_i(\tau,x)|\di{\tau}.
\end{equation}
%
Now, the weak form of the reaction diffusion problems governing the concentrations $u_i^{(j)}$ is given by
%
\begin{multline*}
\int_\Omega\partial_t\bar{u}_i\varphi\di{x}+\int_\Omega D_i(\sigma^{(1)})\nabla\bar{u}_i\cdot\nabla\varphi\di{x}
+\int_\Omega\left(D_i(\sigma^{(1)})-D_i(\sigma^{(2)})\right)\nabla u_i^{(2)}\cdot\nabla\varphi\di{x}\\
=\int_\Omega (R_i(u^{(1)})-R_i(u^{(2)}))\varphi\di{x}
+\int_\Omega\left[A(\sigma^1)(a_iu_i^{(1)}- b_iv^{(1)})-A(\sigma^{(2)})(a_iu_i^{(2)}- b_iv^{(2)})\right]\varphi\di{x}.
\end{multline*}
%
Choosing $\varphi=\bar{u}_i$, integrating over $(0,t)$, and using the uniform positivity of $D_i$, we estimate
%
\begin{multline}\label{eq:initial_estimate}
\frac{1}{2}\|\bar{u}_i(t)\|_{L^2(\Omega)}^2+c_D\|\nabla\bar{u}_i\|_{L^2(S_T\times\Omega)}^2\\
\leq\underbrace{\int_0^t\int_\Omega\left|\left(D_i(\sigma^1)-D_i(\sigma^2)\right)\nabla u_i^{(2)}\cdot\nabla\bar{u}_i\right|\di{x}\di{\tau}}_{I_1}
+\underbrace{\int_0^t\int_\Omega (R_i(u^1)-R_i(u^2))\bar{u}_i\di{x}\di{\tau}}_{I_2}\\
+\underbrace{\int_0^t\int_\Omega\left[A(\sigma^1)-A(\sigma^2)\right](a_iu_i^{(1)}- b_iv^{(1)})\bar{u}_i\di{x}\di{\tau}}_{I_3}
+\underbrace{\int_0^t\int_\Omega A(\sigma^2)(a_i\bar{u}_i- b_i\bar{v})\bar{u}_i\di{x}\di{\tau}}_{I_4}.
\end{multline}
%
For the $I_1$-term, we use \cref{lemma:lipschitz.diffusivity}, \cref{eq:bar_sigma_estimate}, and Fubini to estimate via Hölder's inequality:
%
\[
I_1
\leq L_D\alpha\sum_{k=1}^N a_k\int_0^t\|\bar{u}_k\|_{L^1(0,\tau;L^{q}(\Omega))}\|\nabla u_i^{(2)}\|_{L^{p}(\Omega)}\|\nabla\bar{u}_i\|_{L^2(\Omega)}\di{\tau}
\]
%
with $q=\frac{2p}{p-2}$ and $p>d$ as in the assumption of the theorem.
We note that, with this choice, $q< q^*:=\frac{2d}{d-2}$ (with $q^*=\infty$ when $d=2$) and utilize the continuous and compact embedding from $H^1(\Omega)$ into $L^q(\Omega)$ for all $q<q^*$.
We now use Ehrling's lemma (note that $H^1(\Omega)$ is compactly embedded into $L^q(\Omega)$ since $q<q*$): for all $\delta>0$ there is $C_{\delta}>0$ such that
%
\[
\|u\|_{L^q(\Omega)}\leq C_{\delta}\|u\|_{L^2(\Omega)}+\frac{\delta}{2} \|\nabla u\|_{L^2(\Omega)}\quad  (u \in H^1(\Omega)).
\]
%
This implies
%
\[
I_1
\leq L_D\alpha\sum_{k=1}^N a_k\int_0^t\left(C_{\delta}\|\bar{u}_k\|_{L^1(0,\tau;L^{2}(\Omega))}+\frac{\delta}{2}\|\nabla\bar{u}_k\|_{L^1(0,\tau;L^{2}(\Omega))}\right)\|\nabla u_i^{(2)}\|_{L^{p}(\Omega)}\|\nabla\bar{u}_i\|_{L^2(\Omega)}\di{\tau}.
\]
%
Applying Hölder's inequality once more, we further estimate
%
\begin{equation}\label{eq:estimate_I1}
I_1
\leq L_D\alpha\sqrt{t}\|u_i^{(2)}\|_{L^2((0,t);W^{1,p}(\Omega))}\|\nabla \bar{u}_i\|_{L^2((0,t)\times\Omega)}\sum_{k=1}^N a_k\left(C_\delta\|\bar{u}_k\|_{L^{2}((0,t)\times\Omega)}+\delta\|\nabla\bar{u}_k\|_{L^{2}((0,t)\times\Omega))}\right).
\end{equation}
%
% By assumption, $\nabla u^{(2)}\in L^2(S_T;L^p(\Omega))$.
% Setting $C_1=L_D\alpha\sqrt{T}\max_{k=1,...,N}\left\{a_k\|\nabla u_k^{(2)}\|_{L^2(S_T;L^{p}(\Omega))}\right\}$, we estimate
% %
% \[
% I_1
% \leq C_1\sum_{k=1}^N\left(\left(C_{\delta}+(2\delta)^{-1}\right)\|\bar{u}_k\|^2_{L^{2}(S_T\times\Omega)}+\delta\|\nabla\bar{u}_k\|^2_{L^{2}(S_T\times\Omega)}\right).
% \]
%
The reaction terms $R_i$ are locally Lipschitz continuous, $u_i^{(1)}$ and $u_i^{(2)}$ are essentially bounded (\cref{thm.existence}), and $A$ is uniformly bounded:
%
\begin{equation*}
I_2+I_4
\leq (L_R+Ca_i)\|\bar{u}_i\|_{L^2((0,t)\times\Omega)}^2+C b_i\int_0^t\int_\Omega|\bar{v}||\bar{u}_i|\di{x}\di{\tau}.
\end{equation*}
%
The second term on the right-hand side can be estimated the same as $I_1$ considering that $\bar{\sigma}=\alpha\bar{v}$, which leads to
%
\begin{equation}\label{eq:estimate_I2I4}
I_2+I_4
\leq (L_R+C a_i)\|\bar{u}_i\|_{L^2((0,t)\times\Omega)}^2+C b_i\sqrt{t}\|\bar{u}_i\|_{L^2((0,t)\times\Omega)}\sum_{k=1}^N a_k\|\bar{u}_k\|_{L^2((0,t)\times\Omega)}.
\end{equation}
%
For the remaining term, $I_3$, we use the Lipschitz continuity of $A$ (\cref{lemma:lipschitz_coefficient}) and \cref{eq:bar_sigma_estimate}:
%
\begin{equation}\label{eq:estimate_I3}
I_3
\leq L_A\alpha\sqrt{t}\left(a_i\|u_i^{(1)}\|_{L^\infty((0,t)\times\Omega))}+ b_i\|v^{(1)}\|_{L^\infty((0,t)\times\Omega)}\right)\|\bar{u}_i\|_{L^2((0,t)\times\Omega)}\sum_{k=1}^N a_k\|\bar{u}_k\|_{L^2((0,t)\times\Omega)}
\end{equation}
%
% We set
% %
% \[
% C_1=\max_{i=1,..,N}\left\{\|u_i^{(2)}\|_{L^2(S_T;W^{1,p}(\Omega))},\left(a_i\|u_i^{(1)}\|_{L^2(S_T;L^{4}(\Omega))}+ b_i\|v^{(1)}\|_{L^2(S_T;L^{4}(\Omega))}\right)\right\}
% \]
% and note that $u^{(2)},v^{(2)}\in L^2(S_T;L^4(\Omega))$ via Sobolev-embeddings and $u_i^{(2)}\in L^2(S_T;W^{1,p}(\Omega))$ by assumption.
Using a generic constant $C$, we thereby get
%
% \begin{multline*}
% \frac{1}{2}\|\bar{u}_i(t)\|_{L^2(\Omega)}^2+c_D\|\nabla\bar{u}_i\|_{L^2(S_T\times\Omega)}^2\\
% \leq (L_R+C a_i)\|\bar{u}_i\|_{L^2(S_T\times\Omega)}^2
% \\
% +L_D\alpha\sqrt{t}C_1\|\nabla\bar{u}_i\|_{L^2(S_T\times\Omega)}\sum_{k=1}^N a_k\left(C_\delta\|\bar{u}_k\|_{L^{2}(S_T\times\Omega)}+\delta\|\nabla\bar{u}_k\|_{L^{2}(S_T\times\Omega))}\right)\\
% +\sqrt{t}\left(C  b_i\|\bar{u}_i\|_{L^2(S_T\times\Omega)}+L_V\alpha C_1\|\bar{u}_i\|_{L^2(S_T;L^{4}(\Omega))}\right)\sum_{k=1}^N a_k\|\bar{u}_k\|_{L^2(S_T\times\Omega)}.
% \end{multline*}
%
\begin{multline*}
\|\bar{u}_i(t)\|_{L^2(\Omega)}^2+\|\nabla\bar{u}_i\|_{L^2((0,t)\times\Omega)}^2\\
\leq C\|\bar{u}_i\|_{L^2((0,t)\times\Omega)}^2
+C\bigg(\|\nabla\bar{u}_i\|_{L^2((0,t)\times\Omega)}\sum_{k=1}^N\left(C_\delta\|\bar{u}_k\|_{L^{2}((0,t)\times\Omega)}+\delta\|\nabla\bar{u}_k\|_{L^{2}((0,t)\times\Omega))}\right)\\
+\|\bar{u}_i\|_{L^2((0,t)\times\Omega)}\sum_{k=1}^N\|\bar{u}_k\|_{L^2((0,t)\times\Omega)}\bigg).
\end{multline*}
%
With Cauchy-Schwarz's and Young's inequalities, we can simplify this further 
%
\begin{equation*}
\|\bar{u}_i(t)\|_{L^2(\Omega)}^2+\|\nabla\bar{u}_i\|_{L^2((0,t)\times\Omega)}^2
\leq C\left(\|\bar{u}_i\|_{L^2((0,t)\times\Omega)}^2
+NC_\delta\|\bar{u}\|^2_{L^{2}((0,t)\times\Omega)^N}+N\delta\|\nabla\bar{u}\|^2_{L^{2}((0,t)\times\Omega)^N}\right).
\end{equation*}
%
Now, summing over $i=1,...,N$, taking the supremum over $S_T$, and subsuming the gradient term on the left hand side by choosing $\delta$ to be sufficiently small, we arrive at
%
\[
\|\bar{u}(t)\|_{L^2(\Omega)^N}^2+\|\nabla\bar{u}\|_{L^2(S_T\times\Omega)^N}^2\\
\leq C\|\bar{u}\|^2_{L^2(S_T\times\Omega)^N}
\]
%
which, via Grönwall's inequality, implies $\bar{u}=0$ almost everywhere in $S_T\times\Omega$.
\end{proof}

\end{document}
